\documentclass{report}
\usepackage{amsmath,amssymb}
\setlength{\parindent}{0mm}
\setlength{\parskip}{1em}
\begin{document}
\begin{center}
\rule{6in}{1pt} \
{\large Jeffrey J. Heys \\
{\bf Comparison of Continuous and Discontinuous Galerkin Finite Element Methods for Parabolic Differential Equations}}

Chemical and Biological Engineering \\ Montana State University \\ Box 173920 \\ Bozeman \\ MT 59715
\\
{\tt jeff.heys@gmail.com}\\
Prathish Kumar Rajaraman\\
Garret Vo\\
Glen Hansen\end{center}

A number of different discretization techniques and algorithms have been
developed for approximating the solution of parabolic partial
differential equations. A standard approach, especially for applications
that involve complex geometries, is the classic continuous Galerkin
Finite Element Method (GFEM). This approach has a strong theoretical
foundation and has been widely and successfully applied to this catagory
of differential equations. One challenging catagory of problems, however,
are equations that include an advective term that large relative to the
second-order, diffusive term. For these advection dominated problems, the
continuous GFEM discretization can become unstable and yield highly
inaccurate results. An alternative to the continuous GFEM is the
discontinuous GFEM, and, through the use of a numerical flux term used in
deriving the weak form, the discontinuous approach has the potential to
be much more stable in highly advective problems. However, the
discontinuous GFEM also has significantly more degrees-of-freedom due to
the replication of nodes along element edges and vertices. This paper
compares the computational cost, stability, and accuracy (when possible)
of continuous and discontinuous GFEM for four different test problems
including the advection-diffusion equation, viscous Burgers' equation,
and the Turing pattern formation equation system. The comparison is
performed using as much shared code as possible between the two
algorithms and direct, iterative, and multilevel linear solvers. The
results show that, for implicit time stepping, the continuous GFEM is
typically 5-20 times less computationally expensive than the
discontinuous GFEM using the same finite element mesh and element order.
However, the discontinuous GFEM is significantly more stable than the
continuous GFEM for advection dominated problems and is able to obtain
accurate approximate solutions for cases where the classic, unstabilized
continuous GFEM fails.


\end{document}
